\chapter{Fazit}
\label{ch:fazit}

\section{Beantwortung der Forschungsfrage}
\label{s:fazit-forschungsfrage}

Die empirischen Ergebnisse zeigen, dass der Leistungsunterschied zwischen React~19 und Svelte~5 unter kombinierter Dauerlast nicht rein quantitativ, sondern architektonisch bedingt ist. Entscheidend ist dabei, dass dieselbe Architektureigenschaft jeweils Vor- und Nachteil zugleich erzeugt: Der Reconciliation-Overhead der Fiber-Engine, der bei hoher UI-Dichte zu erhöhten Scripting-Kosten, höherem Speicherbedarf und einem früheren Tipping Point führt, ermöglicht bei moderater Last und hoher Update-Frequenz ein Scheduling-Modell, das den INP stabilisiert. Umgekehrt schließt Sveltes feingranulare Reaktivität, die den Effizienzvorteil bei hoher Dichte begründet, ein vergleichbares Priorisierungsmodell aus und führt unter Extremlast zu einer Umkehr der Long-Task-Rangfolge. Eine vollständige Übersicht der Messbefunde und Hypothesenbewertungen findet sich in \cref{tab:hypotheses-summary}. Kein Framework ist damit uneingeschränkt überlegen; welches besser skaliert, lässt sich nur lastabhängig beantworten. Die vorliegende Untersuchung schließt damit die in \cref{s:ausgangssituation,s:research-gap} identifizierten Lücken der bisherigen Forschung: Nach aktuellem Stand der Literatur lagen bislang weder vergleichbare empirische Messdaten für React~19 mit aktiviertem Compiler und Svelte~5 mit Runes unter kombinierter Dauerlast vor, noch wurden INP und Frame-Budget-Überschreitung als Metriken für diese Frameworks systematisch erfasst und ausgewertet.

\section{Handlungsempfehlung für die Praxis}
\label{s:fazit-empfehlung}

Für datenintensive Echtzeit-Anwendungen wie Monitoring-Dashboards (\cref{s:motivation}) lassen sich auf Basis der empirischen Ergebnisse vier Szenarien unterscheiden:

\noindent Für Anwendungen mit \textbf{hoher UI-Dichte} bietet Svelte~5 strukturelle Vorteile: Der geringere Scripting-Aufwand verschiebt den Tipping Point zu höheren $n$-Werten (\cref{s:results-h2,s:results-h4}); in der Mehrzahl der gemessenen Szenarien erzielt Svelte zudem niedrigere INP-Werte (\cref{ss:results-h1-verdict}). Zwei Ausnahmen schränken diesen Vorteil bei maximaler UI-Dichte ein (\cref{ss:results-h1-phases}).

\noindent In \textbf{speicherkritischen Umgebungen}, etwa auf ressourcenbeschränkten Endgeräten oder bei Anwendungen mit vielen gleichzeitig aktiven Komponenteninstanzen, ist Svelte~5 ebenfalls klar im Vorteil, da der Heap-Verbrauch (wie die internen Datenstrukturen beider Frameworks belegen) strukturell niedriger ausfällt (\cref{s:results-h3}).

\noindent Bei \textbf{moderater UI-Dichte und hoher Update-Frequenz} kann React~19 die bessere Wahl sein: Das prioritätsbasierte Scheduling der Fiber-Engine erzielt in diesem Lastbereich niedrigere INP-Werte (\cref{ss:results-h1-phases}). Dieser Vorteil tritt ausschließlich unterhalb des Tipping Points auf, ab dem Reacts höhere Scripting-Kosten das Frame-Budget überschreiten (\cref{ss:results-h4-tipping}); bei steigender UI-Dichte überwiegen die strukturellen Vorteile von Svelte.

\noindent Bei \textbf{niedriger UI-Dichte} sind die gemessenen Unterschiede insgesamt gering; Svelte~5 zeigt vereinzelt leichte INP-Vorteile (\cref{ss:results-h1-verdict}). Die Framework-Entscheidung kann in diesem Fall anhand anderer Kriterien wie Ökosystem, Entwicklererfahrung oder bestehender Codebasis getroffen werden.

\noindent In allen Szenarien gilt: Die hier ermittelten konkreten Last-Schwellenwerte sind V8- und Chromium-spezifisch; andere JavaScript-Engines können ein abweichendes Skalierungsverhalten zeigen (\cref{s:limitations}).

\section{Ausblick}
\label{s:fazit-ausblick}

Aus den Ergebnissen dieser Arbeit ergeben sich vier Richtungen für weiterführende Forschung.

\noindent\textbf{Übertragbarkeit auf andere JavaScript-Engines.} Die ermittelten Schwellenwerte und der Scripting-Schnittpunkt bei $f \approx 62\,\text{Hz}$ beruhen auf V8s JIT-Optimierungsprofil (\cref{s:limitations}). Eine vergleichende Studie unter SpiderMonkey und JavaScriptCore würde klären, ob die architektonischen Unterschiede zwischen signalbasierter Reaktivität und VDOM-Reconciliation engine-unabhängige Skalierungseigenschaften erzeugen oder ob die beobachteten Tipping-Point-Werte V8-spezifisch sind.

\noindent\textbf{Messung des Compiler-Anteils.} Die vorliegende Arbeit erfasst den kombinierten Effekt aus Fiber-Reconciliation und Compiler-Memoisierung, kann jedoch nicht bestimmen, wie viel jeder Anteil beiträgt (\cref{s:limitations}). Ein direkter A/B-Vergleich mit und ohne Compiler würde diese Trennung ermöglichen und erstmals den isolierten Overhead der automatischen Memoisierung messbar machen. Diese Information ist für die praktische Beurteilung des React Compilers relevant.

\noindent\textbf{Ursachenanalyse der \textit{presentationDelay}-Ausreißer.} Bei $n = 10\,000$ konnten zwei Ausnahmen mit den erhobenen Metriken nicht erklärt werden: Sveltes INP-Nachteil bei $f = 20\,\text{Hz}$ und Reacts \textit{presentationDelay}-Spitze bei $f \approx 30\,\text{Hz}$ (\cref{ss:results-h1-phases,s:limitations}). Eine Chrome-Tracing-basierte Analyse würde die Zerlegung in Layout-, Paint- und Composite-Phasen ermöglichen und damit möglicherweise einen bisher unbekannten Browser-Scheduling-Effekt sichtbar machen.

\noindent\textbf{Übertragbarkeit auf andere Anwendungsmuster.} Der Benchmark simuliert ein teilweise aktualisiertes Grid (\cref{ss:scenario-a}). Ob virtualisierte Listen, reaktive Formulare oder interaktive Diagramme ein vergleichbares Skalierungsverhalten zeigen, bleibt eine offene empirische Frage, insbesondere ob der hier beobachtete Scheduling-Vorteil Reacts bei moderater Last auf anzeigeorientierte oder interaktionsbasierte Anwendungsmuster übertragbar ist.
